\section{A triple-cover obstruction for large ordinary alternating tensor factors}
\label{sec:factorwise-triples}

\paragraph{Status and exact interface.}
These are new root statements HD-FT01--02, question Q2da.
The complete author/root and separate \texttt{/root/imprimitive}
whole-frozen-proof readings pass. The complete detailed independent
source audit is retained; no finite supplemental computation is
required by or claimed as completed in this integration. They concern an actual factorwise tensor decomposition, a
case explicitly outside D-BI19. They do not assume that a projected
factor inherits twelve genuine branch marks.

Let an actual connected finite-group cover of the sphere have coefficient
representation \(V\), degree \(e\), and genuine generating product-one
inertia \(T_1,\ldots,T_r\), with
\[
                  \sum_i\operatorname{rank}(1-T_i)=2e+6.
                  \tag{FT.1}
\]
Suppose there is an actual quotient \(G\longrightarrow P\), with
\(P=A_n\) or \(S_n\) in its natural action, and an actual factorization
\[
             V=W\otimes U,\qquad
             T_i=W(x_i)\otimes U_i
             \quad\hbox{up to a scalar},
             \tag{FT.2}
\]
where \(x_i\) is the image of the inertia in \(P\), \(a=\dim W\),
and \(\dim U\ge1\). The tensor partners and all scalars are unrestricted.
Only the existence of the displayed actual quotient and representation
factorization is assumed; an abstract numerical profile is insufficient.

\subsection{HD-FT01: cubic traces, actual triple covers, and no mark-count premise}

\begin{theorem}[HD-FT01]
Assume \(n>3500000\), \(a\ge n^{120}\), and, with \(f(x)\) the number
of fixed natural letters, assume
\[
 \frac{|\chi_W(x)|}{a}\le
 \begin{cases}
       (f(x)/n)^3,&f(x)>0,\\
       2n^{-6},&f(x)=0 .
 \end{cases}
 \tag{FT.3}
\]
Then no source (FT.1)--(FT.2) exists.
The theorem is independent of the number of genuine projected marks,
the degree or field of \(U\), and the original imaginary quadratic field.
\end{theorem}

\begin{proof}
For a finite-order operator \(B\), put
\(\nu(B)=\dim B-\max_z\dim\ker(B-z)\).
Tensor spectral convolution gives
\[
       \frac{\operatorname{rank}(1-T_i)}e
                   \ge\frac{\nu(W(x_i))}{a}.
       \tag{FT.4}
\]
This also allows arbitrary scalar factors on either side of (FT.2).
If \(x_i=1\), we discard its nonnegative contribution below.
Let \(r_0\) denote the number of remaining indices.

First bound \(r_0\) from the actual total budget.
Since \(2n^{-6}\le n^{-3}\), (FT.3) implies the exact cubic
all-target hypothesis in HD-TW02. Its purely local proof therefore gives
\[
 \frac{\nu(W(x))}{a}
 \ge \frac{(n-4)\operatorname{ind}(x)}{n^2}+\frac1n
 \ge \frac{2n-4}{n^2}
 \qquad(x\ne1).
 \tag{FT.5}
\]
That local proof uses neither transitivity of tensor factors nor a
branch-mark count. From (FT.1), (FT.4), and \(e\ge a\ge n^{120}\),
\[
 r_0\le
 \left(2+\frac6e\right)\frac{n^2}{2n-4}
 \le \left(2+\frac6{n^3}\right)\frac{n^2}{2n-4}<2n.
 \tag{FT.6}
\]
For the last strict inequality it suffices that
\(2n^2-8n-6/n>0\), which holds for \(n\ge5\).

Let \(\Omega=\{1,\ldots,n\}^3\), with ordered entries and repetitions
allowed. For \(x\) of order \(q\), Burnside's orbit formula gives
\[
 \operatorname{ind}(x,\Omega)
       =n^3-\frac1q\sum_{j=0}^{q-1}f(x^j)^3.
       \tag{FT.7}
\]
Cyclic Fourier inversion bounds every eigenvalue multiplicity of \(W(x)\)
by the average absolute character value. Using (FT.3), including the
identity power and every fixed-point-free power, gives
\[
 \frac{\nu(W(x))}{a}
    \ge\frac{\operatorname{ind}(x,\Omega)}{n^3}
                  -\frac2{n^6}.
 \tag{FT.8}
\]
The error can be multiplied by the actual fraction of fixed-point-free
powers, but the displayed weaker bound suffices.

The image \(P\) is three-transitive because \(n\ge5\).
Its action on \(\Omega\) has exactly five orbits: the diagonal,
three equality patterns with exactly two distinct letters, and
the ordered triples of distinct letters. These are actual connected
quotient covers of the given \(P\)-cover. Write their genera as
\(\gamma_1,\gamma_2,\gamma_2,\gamma_2,\gamma_3\).
Riemann--Hurwitz on each orbit yields the exact sum
\[
 \sum_i\operatorname{ind}(x_i,\Omega)
       =2(n^3-5)+2(\gamma_1+3\gamma_2+\gamma_3).
       \tag{FT.9}
\]
Entries with \(x_i=1\) contribute zero, so the same equality holds
over the \(r_0\) retained entries.

Consider also the actual quotient cover on unordered three-element
subsets. Its point stabilizer is
\(H_3=P\cap(S_{n-3}\times S_3)\).
The published Monderer--Neftin theorem states that for \(n>N_1\),
every subgroup not containing \(A_{n-1}\) whose fixed curve has
genus at most one must contain a conjugate of \(A_{n-2}\).
Its explicit constant is \(N_1=3500000\), using the classification
of finite simple groups; the classification-free stated alternative
is \(60000000000\).
The exact input is Theorem1.2, with Remark2.8, of
\cite{HD:MondererNeftin2022}.

Our \(H_3\) cannot contain either required alternating subgroup:
\[
 |H_3|\le6(n-3)!<
           \frac{(n-2)!}{2}=|A_{n-2}|
           \qquad(n>14).
 \tag{FT.10}
\]
Thus the unordered-three-subset curve has genus at least two.
The forget-order map from the distinct ordered-triple curve has
degree six, for both \(P=A_n\) and \(S_n\).
For \(A_n\), every permutation of the selected three letters
can be paired with a permutation of matching parity on their
complement, since \(n-3\ge2\). This proves the degree assertion
also in that case.
Riemann--Hurwitz now gives
\[
                  \gamma_3-1\ge6(2-1),\qquad\gamma_3\ge7.
                  \tag{FT.11}
\]
The other genera are nonnegative. Hence (FT.9) is at least
\(2n^3+4\).

Finally sum (FT.4) and (FT.8), and use (FT.6):
\[
 \frac1e\sum_i\operatorname{rank}(1-T_i)
 \ge2+\frac4{n^3}-\frac{2r_0}{n^6}
 >2+\frac4{n^3}-\frac4{n^5}
 >2+\frac3{n^3}.
 \tag{FT.12}
\]
Since \(e\ge n^{120}\), the last strict excess contributes more
than \(3n^{117}>6\) to the moving-rank sum, contradicting (FT.1).
This proves the theorem.
\end{proof}

\paragraph{Source-reading precision.}
The author revised version dated18 December2020 has Theorem1.2 on
page2, Theorems2.5--2.6 and Remark2.8 on pages7--8, and the complete
function-field proof, Theorem5.1, on pages16--17.
Root read those complete invoked passages, including the full
intermediate Proposition3.1--3.3 and Corollary3.4 proofs.
Publisher metadata verifies the published article, its revision date
and its2022 pagination. The publisher PDF request led to the
subscription preview, so no reading of inaccessible typeset pages
is claimed. Remark2.8 states that its constants come from inspection
of the Neftin--Zieve argument; those numerical constants have not
been independently recomputed in this run.
We invoke the published Monderer--Neftin result with that exact
source scope, not an invented effective version of a merely
asymptotic statement.

\subsection{HD-FT02: an effective growing-tail exclusion with arbitrary tensor partners}

\begin{corollary}[HD-FT02]
In (FT.1)--(FT.2), suppose \(W\) is an irreducible \(S_n\)-module
or its irreducible ordinary restriction to \(A_n\).
If \(n>3500000\) and \(a\ge n^{120}\), the source is impossible.
In particular set
\[
                         N_{\rm tail}=241^{482}.
                         \tag{FT.13}
\]
For \(n\ge N_{\rm tail}\), a surviving such factor must have a
partition label \(\lambda\) satisfying
\[
          \min\{n-\lambda_1,n-\lambda'_1\}\le240 .
          \tag{FT.14}
\]
This assertion allows arbitrary factorwise tensor partners.
It does not claim that these remaining low-tail labels occur.
\end{corollary}

\begin{proof}
The published effective MSP consequence was proved in HD-TW02:
when \(n\ge4000\) and \(a\ge n^{120}\), the absolute character ratio
is at most \((f/n)^3\) for \(f>0\).
Its separate fixed-point-free calculation gives the stronger bound
\[
             |\chi_W(x)|/a<a^{-1/20}\le n^{-6}
                         \qquad(f(x)=0).
                         \tag{FT.15}
\]
Restriction to the ordinary irreducible \(A_n\)-module preserves
both degree and trace. Thus (FT.3) holds, and HD-FT01 applies.
No claim is made here for a nonextendible constituent obtained by
splitting a self-conjugate \(S_n\)-label.

It remains to prove (FT.14) without assuming the earlier long-row
source reduction. Suppose both tails are at least241.
Transpose \(\lambda\) if necessary for the dimension calculation,
and put \(m=\max\{\lambda_1,\lambda'_1\}\).
Since the diagram fits in its first-row by first-column rectangle,
\(m\ge\sqrt n\).
There are at least241 boxes below its first row.
Remove corners from this lower partition until a partition
\(\mu\) of241 remains. Then \((m,\mu)\) is a subpartition of
\(\lambda\). The ordinary branching rule gives
\[
                 f^\lambda\ge f^{(m,\mu)}.
                 \tag{FT.16}
\]
Equivalently, restrict the irreducible successively along a chain
of corner removals; this constituent has positive multiplicity.
The exact published MSP Lemma8(ii), applied to \((m,\mu)\), gives
\[
 f^\lambda\ge\binom{m}{241}f^\mu
       \ge(m/241)^{241}
       \ge\frac{n^{120+1/2}}{241^{241}}
       \ge n^{120}
                   \qquad(n\ge241^{482}).
 \tag{FT.17}
\]
Here \(m\ge241\), so the binomial product bound is valid.
Transposition preserves dimension. Also
\(241^{482}>3500000\).
The first assertion therefore excludes this case and proves
(FT.14).
\end{proof}

\paragraph{Geometric and programme scope.}
This is a new necessary gate for actual factorwise tensor-decomposable
curve suppliers, using the genera of actual finite quotient curves and
the weight-one Chevalley--Weil budget. It neither assumes that a factor
is itself a rank-nine source nor transfers twelve marks to it.
The bound still leaves low tails in unbounded natural degree, other
projective factor types, and at this stage nonextendible split
alternating constituents.
It supplies no cycle and no detector-degree lower bound.
For condition A (cote16,scan050; SIGN_NEXT SN-C03/C04) it removes
a possible supplier range. Condition I (scan049;cote44,Theorem4),
the Cayley--Hamilton projector of scans048--049, the existing
SGA/EGA/FGA spreading and realization interfaces, and the GHOSTS B03
interpretation remain as already stated. Markman, Severn and032
are not premises of this argument.
